\chapter*{Introducere} 
\addcontentsline{toc}{chapter}{Introducere}

În ultimii ani, domeniul procesării semnalelor audio (Audio Signal Processing) se bucură de o popularitate crescută 
datorită multitudinii de aplicații practice existente precum anularea zgomotului (noise cancelling), 
procesarea vorbirii (speech processing), clasificarea sunetelor, dar în special aplicații în domeniul muzical 
(Music Information Retrieval) începând cu clasificarea genurilor muzicale sau a artiștilor pentru crearea de sisteme 
de recomandare sau identificarea încălcării drepturilor de autor prin recunoașterea segmentelor muzicale în videoclipuri, 
pâna la generarea unor segmente muzicale noi pornind de la anumite genuri sau artiști \footnote{\url{https://openai.com/blog/jukebox/}}.

În ceea ce privește problema clasificării genurilor muzicale, progresele sunt semnificative, 
luând în calcul un anumit grad de 
subiectivitate legat de apartenența melodiilor la anumite genuri \cite{subjectivity}. 
Odată cu evoluția tehnologiei, dimensiunea datelor de intrare pentru algoritmii de clasificare nu mai este limitată la 
reprezentări de dimensiune redusă. Dacă la început experimentele se axau pe identificarea caracteristicilor de mici 
dimensiuni dar cu o putere discriminatoare mare prezente în segmentele muzicale (zero crossing rate, MFCC mean, 
spectral flow mean\cite{featureselection}) și aplicarea unor algoritmi precum KMeans, Knn, masini cu vectori suport 
sau rețele neuronale cu un număr redus de straturi, 
acum se utilizează date de dimensiuni mari, chiar neprelucrate (forme de undă) ca date de intrare pentru rețele neuronale, 
de asemenea, de dimensiuni mari, care reușesc să extragă cu succes caracteristicile specifice ale fiecărui gen și 
să încadreze segmentele într-un gen pe baza acestor caracteristici \cite{resut1,resut2}.

În ambele abordări descrise, partea instrumentală a segmentului muzical joacă un rol important, 
în special elementele de percuție care sunt bine sintetizate în reprezentări, 
fie ele de dimensiuni mici (zero crossing rate) sau de dimensiuni mari (spectograme). 
În acest sens, există multe experimente care încearcă să rezolve problema de clasificare utilizând doar 
informații legate de partea instrumentală a segmentelor muzicale\cite{instrumental}. 
În mod surprinzător, puține lucrări s-au concentrat doar pe aspectul vocal în aceasta problemă de clasificare, 
cele existente utilizând doar reprezentări de dimensiuni mici și masini cu vectori suport pentru clasificare\cite{vocal} sau
pe clasificarea artistilor, nu a genurilor muzicale\cite{artists}.

Această lucrare abordează problema clasificării genurilor muzicale, 
încadrând segmente muzicale într-un anumit gen, 
utilizând doar elementul vocal. Lucrarea propune un clasificator  ce combină 
o rețea neuronală convoluțională cu o rețea 
neuronală recurentă și are ca date de intrare spectograme Mel. 

Aplicarea practică a lucrării este un sistem de recomandare sub forma unei aplicații web care,
 pe baza unui input vocal,
 încadrează vocea într-o serie de genuri pe baza rezultatelor clasificării segmentelor
 în care a fost împărțit inputul și sugerează melodii similare existente în setul de date.